% !TEX root = ../../main.tex
% C1.1 - Motivation (translated from the Vietnamese edition).
\section{Motivation}
\label{sec:1.1}

Surveillance cameras are being deployed more and more widely in Vietnam. According to figures from the General Department of Customs cited by the Ministry of Information and Communications, more than 16 million surveillance cameras were imported into Vietnam over the last five years, about 3.2 million a year on average, and the country is estimated to have more than 20 million surveillance cameras in use by 2025~\cite{vnexpress2024camera}. The large number of cameras and their continuous recording produce a very large volume of video.

The value of a camera system does not lie only in deterrence. A review of forty years of research found that surveillance cameras reduce crime by a modest amount~\cite{piza2019cctv}, but their more prominent role is to provide evidence for tracing events after they happen. A study of 251,195 incidents recorded by the British Transport Police found that camera footage was useful to the investigation in about 29\% of cases, and that when useful footage was available, the share of cases solved rose from about 20\% to about 50\%~\cite{ashby2017cctv}. The practical value of camera data therefore depends heavily on whether the right footage can be found.

Many organisations record continuously and need to find when and where a specific person appeared, but reviewing every video by hand is slow and error-prone. When cameras span several areas, searching is usually assigned by area and results are kept in case files for managers, which adds needs for access control and result management.

The initial information about the person is also often incomplete and comes in different forms: a reference image, a verbal description, or a few features such as gender, clothing type and colour, or carried items. A useful system should accept all these descriptions in one matching mechanism rather than a single form of query.

These practical needs motivate the project ``Developing a Person Search System from Multimodal Descriptions''.
